%-------------------------
% Resume in Latex
% Author : Jake Gutierrez
% Based off of: https://github.com/sb2nov/resume
% License : MIT
%------------------------

\documentclass[letterpaper,11pt]{article}

\usepackage{latexsym}
\usepackage[empty]{fullpage}
\usepackage{titlesec}
\usepackage{marvosym}
\usepackage[usenames,dvipsnames]{color}
\usepackage{verbatim}
\usepackage{enumitem}
\usepackage[hidelinks]{hyperref}
\usepackage{fancyhdr}
\usepackage[ngerman]{babel}
\usepackage{tabularx}
\usepackage{iftex}
\ifPDFTeX
  \input{glyphtounicode}
\fi


%----------FONT OPTIONS----------
% sans-serif
% \usepackage[sfdefault]{FiraSans}
% \usepackage[sfdefault]{roboto}
% \usepackage[sfdefault]{noto-sans}
% \usepackage[default]{sourcesanspro}

% serif
% \usepackage{CormorantGaramond}
% \usepackage{charter}


\pagestyle{fancy}
\fancyhf{} % clear all header and footer fields
\fancyfoot{}
\renewcommand{\headrulewidth}{0pt}
\renewcommand{\footrulewidth}{0pt}

% Adjust margins
\addtolength{\oddsidemargin}{-0.5in}
\addtolength{\evensidemargin}{-0.5in}
\addtolength{\textwidth}{1in}
\addtolength{\topmargin}{-.5in}
\addtolength{\textheight}{1.0in}

\urlstyle{same}

\raggedbottom
\raggedright
\setlength{\tabcolsep}{0in}

% Sections formatting
\titleformat{\section}{
  \vspace{-4pt}\scshape\raggedright\large
}{}{0em}{}[\color{black}\titlerule \vspace{-5pt}]

% Ensure that generate pdf is machine readable/ATS parsable
\ifPDFTeX
  \pdfgentounicode=1
\fi

%-------------------------
% Custom commands
\newcommand{\resumeItem}[1]{
  \item\small{
    {#1 \vspace{-2pt}}
  }
}

\newcommand{\resumeSubheading}[4]{
  \vspace{-2pt}\item
    \textbf{#1} $|$ #2 \\
    \textit{\small#3} $|$ \textit{\small #4} \\
  \vspace{-7pt}
}

\newcommand{\resumeSubSubheading}[2]{
    \item
    \begin{tabular*}{0.97\textwidth}{l@{\extracolsep{\fill}}r}
      \textit{\small#1} & \textit{\small #2} \\
    \end{tabular*}\vspace{-7pt}
}

\newcommand{\resumeProjectHeading}[2]{
    \item
    \begin{tabular*}{0.97\textwidth}{l@{\extracolsep{\fill}}r}
      \small#1 & #2 \\
    \end{tabular*}\vspace{-7pt}
}

\newcommand{\resumeSubItem}[1]{\resumeItem{#1}\vspace{-4pt}}

\renewcommand\labelitemii{$\vcenter{\hbox{\tiny$\bullet$}}$}

\newcommand{\resumeSubHeadingListStart}{\begin{itemize}[leftmargin=0.15in, label={}]}
\newcommand{\resumeSubHeadingListEnd}{\end{itemize}}
\newcommand{\resumeItemListStart}{\begin{itemize}}
\newcommand{\resumeItemListEnd}{\end{itemize}\vspace{-5pt}}

%-------------------------------------------
%%%%%%  RESUME STARTS HERE  %%%%%%%%%%%%%%%%%%%%%%%%%%%%


\begin{document}

%----------HEADING----------
\begin{center}
    \textbf{\Huge \scshape Erik Weidenauer} \\ \vspace{1pt}
    \large \textit{Full-Stack Entwickler} \\ \vspace{1pt}
    \small \href{mailto:weidenauer.erik@outlook.com}{\underline{weidenauer.erik@outlook.com}} $|$
    \href{https://weidenauererik.com}{\underline{weidenauererik.com}} $|$
    \href{https://linkedin.com/in/erik-weidenauer}{\underline{linkedin.com/in/erik-weidenauer}} $|$
    \href{https://github.com/WeidenauerErik}{\underline{github.com/WeidenauerErik}} \\ \vspace{1pt}
    \small +43 664 75596725 $|$ Wien, \"Osterreich
\end{center}


%-----------EDUCATION-----------
\section{Ausbildung}
  \resumeSubHeadingListStart
    \resumeSubheading
      {HTL Rennweg}{Wien, \"Osterreich}
      {Informationstechnologie / Note: 1,8}{Sep 2021 -- Jun 2026}
      \resumeItemListStart
        \resumeItem{Full-Stack-Webentwicklung: Java / Spring Boot und Symfony (PHP) f\"ur das Backend, Vue.js mit Capacitor f\"ur plattform\"ubergreifende Frontends}
        \resumeItem{Netzwerktechnik \& Cybersecurity: Konfiguration von Cisco-Routern/-Switches, Schutz vor g\"angigen Schwachstellen (z.\,B. SQL-Injection)}
        \resumeItem{Multimedia \& Design: Adobe Creative Suite (Illustrator, Photoshop, InDesign), Fotografie und Filmproduktion}
        \resumeItem{Diplomarbeit: DanceMaster --- plattform\"ubergreifende AR-Tanz-Lernapplikation (siehe Projekte)}
      \resumeItemListEnd
  \resumeSubHeadingListEnd


%-----------EXPERIENCE-----------
\section{Berufserfahrung}
  \resumeSubHeadingListStart

    \resumeSubheading
      {\"Osterreichisches Bundesheer}{Okt 2026 -- Heute}
      {Grundwehrdienst \& Heeresleistungssportzentrum (Voltigieren)}{\"Osterreich}
      \vspace{10pt}

    \resumeSubheading
      {Freiberuflicher Softwareentwickler}{Jun 2026 -- Aug 2026}
      {Selbstst\"andig}{Remote}
      \resumeItemListStart
        \resumeItem{Entwicklung ma{\ss}geschneiderter Softwarel\"osungen f\"ur Unternehmen und Privatkunden, Automatisierung von Prozessen und End-to-End-Lieferung digitaler Produkte}
        \resumeItem{Umsetzung von sauberem, wartbarem Code in Webanwendungen, Desktop-Tools, Backend-Systemen und APIs mit PostgreSQL und Doctrine, abgestimmt auf die Anforderungen jedes Kunden}
      \resumeItemListEnd

    \resumeSubheading
      {Website-Administrator}{Jan 2023 -- Feb 2026}
      {Naturfreunde \"Osterreich}{Remote}
      \resumeItemListStart
        \resumeItem{Eigenverantwortliche Administration und Wartung der offiziellen Website der Ortsgruppe Neubau}
        \resumeItem{Verwaltung von Content-Updates und regelm\"a{\ss}ige Wartung zur Sicherstellung der technischen Stabilit\"at}
        \resumeItem{Optimierung der CMS-basierten Plattform f\"ur eine zuverl\"assige, ansprechende digitale Pr\"asenz der Community}
      \resumeItemListEnd

    \resumeSubheading
      {Praktikant Netzwerktechnik}{Aug 2024 -- Aug 2024}
      {CANCOM Austria}{Wien, \"Osterreich (Hybrid)}
      \resumeItemListStart
        \resumeItem{Fortinet-NSE-3-Zertifizierung erworben und Fachwissen im Bereich Netzwerksicherheit und Infrastruktur vertieft}
        \resumeItem{Entwicklung von Python-Skripten zur Netzwerkautomatisierung zur Verbesserung von Effizienz und Zuverl\"assigkeit}
      \resumeItemListEnd

    \resumeSubheading
      {Praktikant Netzwerktechnik}{Aug 2023 -- Aug 2023}
      {CANCOM Austria}{Wien, \"Osterreich (Hybrid)}
      \resumeItemListStart
        \resumeItem{Konfiguration und Verwaltung von Cisco-Switches sowie Palo-Alto-Networks- und Fortinet-FortiGate-Firewalls in komplexen Netzwerkumgebungen}
        \resumeItem{Fortinet-NSE-1- und NSE-2-Zertifizierungen erworben, fundierte Grundlagen in Cybersecurity und Bedrohungsabwehr nachgewiesen}
      \resumeItemListEnd

  \resumeSubHeadingListEnd


%-----------PROJECTS-----------
\section{Projekte}
    \resumeSubHeadingListStart
      \resumeProjectHeading
          {\textbf{DanceMaster, Diplomarbeit} $|$ \emph{Unity, PostgreSQL, Symfony, JWT, Docker}}{Jul 2025 -- Heute}
          \resumeItemListStart
            \resumeItem{Plattform\"ubergreifende AR-Mobile-App in Unity 6 (iOS/Android) mit AR Foundation, ARKit und ARCore f\"ur Echtzeit-Spatial-Tracking und interaktive Fu{\ss}verfolgungs-Animationen entwickelt}
            \resumeItem{Symfony-REST-API mit JWT-Authentifizierung als Orchestrierungsschicht f\"ur Nutzerdaten und Choreografie-Synchronisation entwickelt; mit Docker containerisiert}
          \resumeItemListEnd
      \resumeProjectHeading
          {\textbf{Heatseekers} $|$ \emph{Python, OpenCV, Flask, Vue.js, SCSS, Bash}}{Okt 2023 -- Feb 2024}
          \resumeItemListStart
            \resumeItem{Python/OpenCV-Hintergrunddienst f\"ur Echtzeit-Bewegungs- und W\"armeerkennung aus Live-Kamerastreams mittels ressourcenschonender Frame-Differenzierung entwickelt}
            \resumeItem{Flask-REST-API und Vue.js/SCSS-Dashboard f\"ur Live-Videostreams und Erkennungsprotokolle entwickelt; Headless-Deployment mit Bash-Skripten automatisiert}
          \resumeItemListEnd
      \resumeProjectHeading
          {\textbf{TaskMaster, Zentraler Klassen-Aufgabenmanager} $|$ \emph{Spring Boot, Vue.js, Thymeleaf}}{Apr 2022 -- Sept 2022}
          \resumeItemListStart
            \resumeItem{Spring-Boot-Backend mit rollenbasierter Zugriffskontrolle zur Trennung von Lehrer- und Sch\"uler-Workflows entwickelt}
            \resumeItem{Vue.js-Frontend mit serverseitigem Thymeleaf-Templating f\"ur dynamisches, sicheres Aufgaben-Tracking kombiniert}
          \resumeItemListEnd
    \resumeSubHeadingListEnd


%-----------CERTIFICATES-----------
\section{Zertifikate}
 \begin{itemize}[leftmargin=0.15in, label={}]
    \small{\item{
     \textbf{Career Essentials in Generative AI} $|$ Microsoft \& LinkedIn, Jul 2026 \\
     \textbf{Career Essentials in GitHub Professional Certificate} $|$ GitHub, Jul 2026 \\
     \textbf{Artificial Intelligence Fundamentals} $|$ IBM, Jul 2026 \\
     \textbf{Introduction to Responsible AI} $|$ Google, Jul 2026 \\
     \textbf{Amazon Translate: Getting Started} $|$ AWS, Feb 2026 \\
     \textbf{Amazon Location Service: Getting Started} $|$ AWS, Feb 2026 \\
     \textbf{SkillDisplay Profile (Symfony, Spring Boot, VueJs, Drupal)} $|$ SkillDisplay GmbH, Jun 2026 \\
     \textbf{Drupal in a Day} $|$ Drupal, Okt 2025 \\
     \textbf{Communication and Team Building Seminar} $|$ HTL Rennweg, M\"arz 2025 \\
     \textbf{NSE 1, NSE 2, NSE 3} $|$ Fortinet, 2023 -- 2024
    }}
 \end{itemize}


%-----------AWARDS-----------
\section{Auszeichnungen}
 \begin{itemize}[leftmargin=0.15in, label={}]
    \small{\item{
     \textbf{Weltmeister 2025} $|$ Voltigieren (Gruppe) $|$ FEI, Aug 2025 \\
     \textbf{8. Platz Weltmeisterschaft 2025} $|$ Voltigieren (Einzel) $|$ FEI, Jul 2025 \\
     \textbf{8. Platz Europameisterschaft 2024} $|$ Voltigieren (Einzel) $|$ FEI, Jul 2024 \\
     \textbf{12. Platz Weltmeisterschaft 2023} $|$ Voltigieren (Einzel) $|$ FEI, Jul 2023
    }}
 \end{itemize}


%-----------PROGRAMMING SKILLS-----------
\section{Technische F\"ahigkeiten}
 \begin{itemize}[leftmargin=0.15in, label={}]
    \small{\item{
     \textbf{Programmierung}{: Vue.js, Spring Boot, Symfony, Java, PHP, Python, C} \\
     \textbf{CMS}{: WordPress, Typo3, Drupal} \\
     \textbf{Design \& Modellierung}{: Blender, Illustrator, Photoshop, InDesign} \\
     \textbf{Cybersecurity \& Netzwerktechnik}{: Cisco Networking, Fortinet, Palo Alto Networks} \\
     \textbf{Tools}{: Git / GitHub, Docker / DockerHub, JetBrains-Produkte, VSCode}
    }}
 \end{itemize}


%-----------LANGUAGES-----------
\section{Sprachen}
 \begin{itemize}[leftmargin=0.15in, label={}]
    \small{\item{
     \textbf{Deutsch}{: Muttersprache} \\
     \textbf{Englisch}{: B2 (Matura / Fortgeschritten)}
    }}
 \end{itemize}


%-------------------------------------------
\end{document}
